% Young 表講義第 9 段：附錄 A 組合的各種變形（動畫 F56–F61）
% 數字由 fulton/make_figs.py 的 check_numbers_appAB() 核對（codes/utils/fulton_appendix.py）；\FulEvacSmall、\FulGaleRyser 由程式產生。

\section{組合的各種變形（動畫 F56--F61）}\label{sec:f-appA}

\begin{keybox}[這個附錄在做什麼（先看這裡）]
第 1--5 章的主角是「列插入」與「滑動」。附錄 A 換幾種玩法，看看還能得到什麼：
\begin{enumerate}[nosep]
\item \textbf{把字母表倒過來}：每個字倒著讀、每個字母換成「對面」的字母。表會跟著變成另一張表 $T^*$（evacuation）。
\item \textbf{一行一行地擠}（行插入）：和列插入得到同一個 $P$。
\item \textbf{換一種排序}（Burge 對應）、\textbf{從矩陣的四個角出發}：得到四對表，彼此只差 $^*$。
\item \textbf{0-1 矩陣}：得到形狀互為共軛的一對表；能不能填出列和、行和給定的 0-1 矩陣，由支配序決定（Gale–Ryser）。
\item \textbf{鑰匙}：把表的各行換順序排，最左、最右那一行只和它的長度有關。
\end{enumerate}
這裡的東西都可以直接拿紙筆算，每一小節都至少有一個從頭算到尾的例子。
\end{keybox}

% ------------------------------------------------------------------
\subsection{把字母表倒過來：對偶與 evacuation（F56）}\label{sec:f-evac}

\textbf{對偶字母}（\fulref{p.183}）。字母表 $[N]=\{1<2<\cdots<N\}$ 的對偶是把順序倒過來：$x^*=N+1-x$，所以 $x<y\iff x^*>y^*$。
一個字 $w=x_1x_2\cdots x_r$ 的對偶字是\textbf{倒過來讀、再換成對偶}：
\[
w^*=x_r^*\cdots x_2^*x_1^* .
\]
\textbf{例 A.1}（$N=3$）。$w=3121$：倒過來是 $1213$，再把 $1,2,3$ 換成 $3,2,1$，得 $w^*=3231$。

\textbf{Knuth 關係保持不變}：$xzy\equiv zxy$（$x<y<z$）變成 $y^*z^*x^*\equiv y^*x^*z^*$，而 $z^*<y^*<x^*$，這又是一條 Knuth 關係。所以等價的字，對偶以後還是等價。

\textbf{evacuation：拆積木}（\fulref{p.184}）。給一張表 $T$：
\begin{enumerate}[nosep]
\item 拿走左上角的數 $x$；
\item 剩下的是一張斜表，用第 1 章的滑動把它滑成一般的表 $\Delta T$——它少了一格；
\item 在那個少掉的格子填上 $x^*$；
\item 對 $\Delta T$ 重複，直到全部拿完。填好的表就是 $T^*$。
\end{enumerate}

\textbf{例 A.2}（$T=\yts{12,3}$，$N=3$，逐步做）。拿走 1，剩下 $\yts{{}2,3}$ 滑成 $\yts{2,3}$，空出第 1 列第 2 格，填 $1^*=3$；
拿走 2，剩下的滑成 $\yts{3}$，空出第 2 列第 1 格，填 $2^*=2$；最後拿走 3，空出左上角，填 $3^*=1$：
\begin{center}
\FulEvacSmall
\end{center}
所以 $T^*=\yts{13,2}$。再做一次會回到 $\yts{12,3}$（對偶定理 (2)）。另一個例子：$\yts{112,23}^*=\yts{123,23}$。

\textbf{例 A.3}（書上的例子，\fulref{p.184}，$N=6$）。$T=\yts{1223,2355,446,56}$。第一步拿走 1，滑成 $\Delta T=\yts{2223,3455,466,5}$，空出第 4 列第 2 格，填 $1^*$。
一路做下去得
\[
T^*=\begin{array}{cccc}6^*&6^*&5^*&4^*\\5^*&5^*&4^*&2^*\\3^*&3^*&2^*\\2^*&1^*\end{array}
\;=\;\yts{1123,2235,445,56}\quad(\text{用 }x^*=7-x).
\]

\begin{thmbox}[對偶定理 \fulref{p.184}]
(1) $T^*$ 與 $T$ 同形狀；(2) $(T^*)^*=T$；(3) $P(w^*)=P(w)^*$，等價地 $(T\cdot U)^*=U^*\cdot T^*$；
(4) 若陣列 $\binom{u_1\cdots u_r}{v_1\cdots v_r}$ 對應 $(P,Q)$，則 $\binom{u_r^*\cdots u_1^*}{v_r^*\cdots v_1^*}$ 對應 $(P^*,Q^*)$。
\end{thmbox}
\textbf{例 A.4}（驗證 (3)）。$P(3121)=\yts{11,2,3}$，$P(w^*)=P(3231)=\yts{13,2,3}$，而 $\yts{11,2,3}^*=\yts{13,2,3}$ $\checkmark$。
\textbf{為什麼形狀不變}：第 3 章的 $L(w,k)$（$k$ 條互不相交的弱遞增子序列的最大總長）對 $w$ 與 $w^*$ 一樣——把子序列倒著讀就是了。

\textbf{一個應用}（\fulref{p.185}）。把 $[T,U]\mapsto[U^*,T^*]$，就得到第 5 章兩組計數之間的一一對應，所以 $c^\nu_{\lambda\mu}=c^\nu_{\mu\lambda}$。
§A.3 把這種「同一組字、不同的切法」做得更徹底：研究兩個斜表何時「在任何滑動下都以同樣方式改變形狀」，
由此在不同形狀的 LR 斜表之間建立一一對應（形狀改變定理，\fulref{p.191}）。這一節初讀可以跳過。

% ------------------------------------------------------------------
\subsection{一行一行地擠：行插入（F57）}\label{sec:f-colbump}

\textbf{規則}（\fulref{p.186}）。把 $x$ 行插入表 $T$（記作 $x\to T$）：看第一行，若 $x$ 比這一行所有數都大，就放在這一行的底部；
否則擠掉這一行裡\textbf{最上面的那個 $\ge x$ 的數}，被擠出來的數去下一行，重複。

\textbf{例 A.5}。
\begin{itemize}[nosep]
\item $1\to\yts{13,2}$：第一行是 $1,2$，最上面的 $\ge1$ 是 1，換成 1、擠出 1；第二行只有 3，$\ge1$ 的是 3，擠出 3；第三行是空的，3 放在頂端：$\yts{113,2}$。
\item $2\to\yts{13,2}$：第一行 $1,2$ 中 $\ge2$ 的最上面是 2，擠出 2；第二行 3 被擠出；第三行放 3：$\yts{123,2}$。
\item $2\to\yts{12,3}$：第一行 $1,3$ 中 $\ge2$ 的是 3，換成 2、擠出 3；第二行只有 2，3 比它大，放在底部：$\yts{12,23}$。
\end{itemize}

\begin{thmbox}[命題 1 \fulref{p.188}（$P$ 的部分）]
字 $w=v_1\cdots v_r$ 的表 $P(w)$，也可以\textbf{倒著}用行插入得到：先放 $v_r$，再行插入 $v_{r-1}$，…，最後 $v_1$。（$Q$ 用「滑進來」得到，同樣是 R-S-K 的 $Q$。）
\end{thmbox}

\textbf{例 A.6}（打擂台）。$w=24135$。左邊用列插入 $2,4,1,3,5$；右邊用行插入 $5,3,1,4,2$。兩邊最後都是
\[
P=\yts{135,24}.
\]
\textbf{例 A.7}（字母都不同時，倒著讀＝轉置）。$P(53142)=\yts{12,34,5}$，正好是 $\yts{135,24}$ 的轉置。
理由：倒過來的字，遞增子序列變成遞減子序列，所以列與行的角色對調（第 3 章）。練習 6（\fulref{p.189}）進一步說 $Q(w^{\rm rev})$ 與 $Q(w^*)$ 互為轉置。

% ------------------------------------------------------------------
\subsection{換一種排序：Burge 對應（F58）}\label{sec:f-burge}

R-S-K（第 4 章）把矩陣寫成\textbf{字典序}的兩列陣列，用列插入。如果把排序換成\textbf{反字典序}——上列由大到小，上列相同時下列由小到大——
再配上行插入，也會得到同形狀的一對表（\S A.4.1，\fulref{p.198}）。

\textbf{例 A.8}（書上的矩陣，\fulref{p.199}）。
\[
A=\begin{pmatrix}1&2&0\\0&1&1\\1&1&1\end{pmatrix}\ \longrightarrow\
\text{成對 }(1,1),(1,2),(1,2),(2,2),(2,3),(3,1),(3,2),(3,3).
\]
反字典序排成 $\binom{3\,3\,3\,2\,2\,1\,1\,1}{1\,2\,3\,2\,3\,1\,2\,2}$。作法 (1b)：從最右邊開始，依序行插入 $v_8,v_7,\dots,v_1$（$2,2,1,3,2,3,2,1$），
每次把對應的 $u$（$1,1,1,2,2,3,3,3$）放進新長出來的格子：
\[
R=\yts{11222,23,3},\qquad S=\yts{11133,22,3}\qquad(\text{形狀 }(5,2,1)).
\]
對照同一個矩陣的 R-S-K：$P=\yts{112223,23}$、$Q=\yts{111223,33}$（形狀 $(6,2)$）。\textbf{同一個矩陣，兩個不同的對應。}
命題 2（\fulref{p.200}）：作法 (2b)（列插入＋滑動）與 (3b)（矩陣球、NW 序）也得到同一對 $(R,S)$。

% ------------------------------------------------------------------
\subsection{從四個角看矩陣（F59）}\label{sec:f-corners}

R-S-K 的矩陣球（第 4 章）是從\textbf{左上角}開始編號的。從另外三個角開始，或者等價地，把陣列的上列或下列換成對偶字母，會怎樣？

\begin{thmbox}[定理 1 \fulref{p.202}（弱序的部分）]
從左上、右下、左下、右上四個角出發，依序得到 $(P,Q)$、$(P^*,Q^*)$、$(R,S^*)$、$(R^*,S)$，其中 $(P,Q)$ 是 R-S-K、$(R,S)$ 是 Burge。
\end{thmbox}

\textbf{例 A.9}（同一個 $A$）。
\begin{center}
\renewcommand{\arraystretch}{2.6}
\begin{tabular}{@{}llll@{}}
\toprule
角 & 做法（等價的陣列） & 結果 & 數值\\
\midrule
左上 nw & $\binom{u}{v}$ 字典序、R-S-K & $(P,Q)$ & $\yts{112223,23}$，$\yts{111223,33}$\\
右下 se & $\binom{u^*}{v^*}$ & $(P^*,Q^*)$ & $\yts{112223,23}$，$\yts{111333,22}$\\
左下 sw & $\binom{u^*}{v}$ & $(R,S^*)$ & $\yts{11222,23,3}$，$\yts{11123,23,3}$\\
右上 ne & $\binom{u}{v^*}$ & $(R^*,S)$ & $\yts{11223,22,3}$，$\yts{11133,22,3}$\\
\bottomrule
\end{tabular}
\end{center}
（這個 $P$ 剛好滿足 $P^*=P$。）程式另外用所有 $2\times2$、元素 $0,1,2$ 的矩陣核對過這個定理。

% ------------------------------------------------------------------
\subsection{填格子遊戲：0-1 矩陣與 Gale--Ryser（F60）}\label{sec:f-zeroone}

\textbf{遊戲}：能不能填出一個 0-1 矩陣，第 $i$ 列的和是 $\lambda_i$、第 $j$ 行的和是 $\mu_j$？
\begin{itemize}[nosep]
\item 列和 $(3,2,1)$、行和 $(2,2,2)$？
\item 列和 $(3,3)$、行和 $(3,1,1,1)$？
\end{itemize}
先自己試試看再往下讀。

\textbf{Knuth 的對應}（命題 3，\fulref{p.203}）。0-1 矩陣的陣列沒有重複的對 $\binom ij$。依序\textbf{行插入} $v_1,v_2,\dots,v_r$，
並把 $u_i$ 放在新格子的\textbf{共軛位置}（第 $a$ 列第 $b$ 格 $\mapsto$ 第 $b$ 列第 $a$ 格）。結果是一對形狀\textbf{互為共軛}的表 $\{P,Q\}$，而且這是一一對應。
\textbf{例 A.10}。
\[
\begin{pmatrix}0&1&1\\1&0&1\\1&0&1\end{pmatrix}\ \longrightarrow\ P=\yts{112,333},\qquad Q=\yts{11,22,33}
\]
（$P$ 的形狀 $(3,3)$，$Q$ 的形狀 $(2,2,2)$，互為共軛。）程式核對過：$3\times3$ 的 511 個非零 0-1 矩陣恰好對應 511 個不同的表對。

\textbf{對偶 Cauchy 恆等式}（Littlewood，\fulref{p.204}）：
\[
\prod_{i,j}(1+x_iy_j)=\sum_\lambda s_\lambda(x_1,\dots,x_n)\,s_{\tilde\lambda}(y_1,\dots,y_m).
\]
（第 4 章的 $\prod(1-x_iy_j)^{-1}=\sum s_\lambda(x)s_\lambda(y)$ 的「對偶版」。）比較係數得練習 19：
\[
\#\{\text{列和 }\lambda\text{、行和 }\mu\text{ 的 0-1 矩陣}\}=\sum_\nu K_{\tilde\nu\lambda}K_{\nu\mu}.
\]

\textbf{揭曉}。列和 $(3,2,1)$、行和 $(2,2,2)$ 恰有 3 個（上面公式也給 3）：
\begin{center}
\FulGaleRyser
\end{center}
列和 $(3,3)$、行和 $(3,1,1,1)$ 一個都沒有：第 1 行要 3 個 1，但只有 2 列。

\begin{thmbox}[Gale--Ryser（練習 20，\fulref{p.204}）]
列和 $\lambda$、行和 $\mu$ 的 0-1 矩陣存在 $\iff$ $\mu\trianglelefteq\tilde\lambda$（$\tilde\lambda$ 是 $\lambda$ 的共軛，$\trianglelefteq$ 是第 2 章的支配序：部分和逐項 $\le$）。
\end{thmbox}
\textbf{驗算}：$\widetilde{(3,2,1)}=(3,2,1)$，$(2,2,2)$ 的部分和 $2,4,6$ 逐項 $\le3,5,6$ $\checkmark$。$\widetilde{(3,3)}=(2,2,2)$，$(3,1,1,1)$ 的第一個部分和 $3>2$，不行。
再兩個：$(2,2)/(2,2)$ 只有全 1 矩陣一個；$(3,1)/(2,1,1)$ 也只有一個（第 1 列全 1）。

% ------------------------------------------------------------------
\subsection{把行重新排隊：鑰匙（F61）}\label{sec:f-keys}

\textbf{問題}：一張表的各行換個順序排（例如短的行放到前面），還能不能做成一張斜表、rectify 回原來的表？

\begin{thmbox}[命題 7 \fulref{p.208}]
給定 $T$ 各行長度的任一種排法，恰有一個斜表（緊湊形式）有這些行長、而且 rectify 成 $T$。這種斜表叫 \textbf{frank}。
\end{thmbox}
\textbf{緊湊形式}：相鄰兩行，左邊較長就頂端對齊、右邊較長就底端對齊。

\textbf{例 A.11}（書上的例子，\fulref{p.209}）。$T=\yts{1122,233,4}$，行長是 $3,2,2,1$。把行長排成 $(2,3,2,1)$，唯一的 frank 斜表各行（由上而下）是
\[
\{1,2\},\quad\{1,3,4\},\quad\{2,3\},\quad\{2\}
\]
（第一行比第二行短，所以底端對齊、往下錯一格）。

\textbf{鑰匙}（推論 1、2，\fulref{p.209--210}）。frank 斜表最左邊那一行的內容，只跟它的長度 $c$ 有關，記作 $\mathcal L_c$；最右邊那一行同理記作 $\mathcal R_c$。
而且 $c<d$ 時 $\mathcal L_c\subset\mathcal L_d$、$\mathcal R_c\subset\mathcal R_d$。把它們排成和 $T$ 同形狀的表，就是\textbf{左鑰匙} $K_-(T)$ 與\textbf{右鑰匙} $K_+(T)$：
\[
T=\yts{1122,233,4}:\qquad K_-(T)=\yts{1111,222,4},\qquad K_+(T)=\yts{2222,333,4}.
\]
\textbf{例 A.12}（小的）。$T=\yts{12,3}$：$K_-(T)=\yts{11,3}$、$K_+(T)=\yts{22,3}$。鑰匙的特徵：每一行的數字集合都包含右邊那一行的集合（推論 2）。
鑰匙是 Lascoux–Schützenberger 用來研究旗流形上線叢截面的標準基底的工具（\S10.5 末提到的「左鑰匙」）。

\begin{notebox}[練習（答案在後面）]
\begin{enumerate}[nosep]
\item $N=4$，求 $w=4132$ 的 $w^*$。
\item 做 $\yts{12,3}$ 的 evacuation 兩次，驗證回到原表。
\item 行插入：$3\to\yts{12,3}$。
\item 列和 $(2,2,1)$、行和 $(3,1,1)$ 的 0-1 矩陣存在嗎？
\end{enumerate}
\textbf{答案}：1. 倒過來 $2314$，換成 $5-x$：$w^*=3241$。
2. $\yts{12,3}\to\yts{13,2}\to\yts{12,3}$。
3. 第一行 $1,3$ 中 $\ge3$ 的是 3，換成 3（不變）、擠出 3；第二行 2 比 3 小，3 放在底部：$\yts{12,33}$。
4. $\widetilde{(2,2,1)}=(3,2)$，$(3,1,1)$ 的部分和 $3,4,5$ 對 $3,5,5$ 逐項 $\le$，存在（例如 $\left(\begin{smallmatrix}1&1&0\\1&0&1\\1&0&0\end{smallmatrix}\right)$）。
\end{notebox}
